\documentclass[11pt,a4paper]{article}
\usepackage[T1]{fontenc}
\usepackage{lmodern}
\usepackage[margin=28mm]{geometry}
\usepackage{amsmath,amssymb,amsthm,booktabs}
\usepackage{microtype}

\usepackage[hidelinks]{hyperref}
\hypersetup{pdftitle={Maximal sorting time under clumped-letter insertion
  in two short multiplicity profiles},pdfauthor={Xiaojian Zeng}}
\newtheorem{theorem}{Theorem}
\newtheorem{lemma}[theorem]{Lemma}
\newtheorem{corollary}[theorem]{Corollary}
\newtheorem{proposition}[theorem]{Proposition}
\DeclareMathOperator{\trunc}{trunc}
\newcommand{\D}{\mathrm{D}}




\title{Maximal sorting time under clumped-letter insertion\\
  in two short multiplicity profiles}
\author{Xiaojian Zeng\\
  \small PBC School of Finance, Tsinghua University, 100190\\
  \small\texttt{xiaojianzeng@foxmail.com}}
\date{}

\begin{document}
\maketitle
\begin{abstract}
Let $T$ denote the number of passes needed to sort a word through the greedy
$aba$-avoiding stack, where sorting means making each letter contiguous.
For $r\geq3$, let $q$ contain $r$ letters, each twice, with at most one letter
allowed a third occurrence. We prove that inserting any finite number of
additional letters, each confined to one contiguous block, preserves the
condition $T(q)=r$. There is no restriction on the lengths of the inserted
blocks, and adjacent equal letters in $q$ are allowed. The proof combines
clumped-letter deletion with induction on $r$ and a complete computer-assisted
base of $28{,}800$ marked three-letter inputs. A consequence is a first-letter
deletion reduction for maximally slow words of length $2k+2$, for $k\geq4$.
The uniform classification and enumeration of that length layer remain
unresolved here.
\end{abstract}

\section{Setting and contribution}

A word represents a set partition by putting positions with equal letters
in the same block. Two words are identified exactly when one is obtained
from the other by a bijective renaming of its alphabet. Position order is
retained. Defant and Kravitz introduced foot-sorting as a stack-sorting
analogue for such objects~\cite{DefantKravitz2024}. We use Xia's deterministic
map~\cite[Definition 2.5]{Xia2024}, denoted here by $\phi$.

Start with an empty stack and output, and read the input from left to right.
Push the next letter if the resulting stack avoids a subsequence $aba$ with
$a\ne b$. Otherwise pop the top letter into the output and retry the same
input letter. Finally drain the stack into the output. The forbidden pattern
need not be consecutive. Either stack orientation gives the same forbidden
triples, since $aba$ is unchanged by reversal.

A letter is \emph{clumped} when all its occurrences form one nonempty
contiguous block. A word is \emph{sorted} if all its letters are clumped;
there is no prescribed order on the blocks. Put
\[
 T(w)=\min\{t\geq0:\phi^t(w)\text{ is sorted}\}.
\]
In particular, sorted words have time zero. A sorted word stays sorted
under $\phi$, which reverses the order of its blocks. Xia proved $T(w)\leq r$
for words on $r$ letters, with sharpness for $r\geq3$, and asked for
the words attaining it~\cite[Theorem 3.1 and Question 6.1]{Xia2024}.

Choi, Gan, Li, and Zhu showed that, for $r\geq3$, the only such word of
length at most $2r$, up to renaming, is $(a_1\cdots a_r)^2$. At length
$2r+1$ they obtained
\[
 \binom{r+1}{2}+2\binom r2
\]
classes~\cite[Theorems 1.1--1.2]{ChoiGanLiZhu2025}. Their Lemma 3.5 also
counts the length-$2r+2$ subclass whose first letter occurs four times and
all others twice. These results are prior work. Related work of Ganesh,
Xia, and Ying concerns image membership and numbers of preimages for the
stack maps~\cite[Sections 3--4]{GaneshXiaYing2024}.

For a set of letters $A$, write $\D_A(w)$ for deletion of all occurrences
of letters in $A$, and abbreviate $\D_{\{a\}}$ to $\D_a$.
The following theorem concerns the two shortest multiplicity profiles in
which every letter is repeated.

\begin{theorem}[Clumped-letter insertion; computer-assisted base]
\label{thm:main}
Let $r\geq3$ and let $q$ use exactly $r$ letters, each occurring twice,
except that at most one may occur three times. Let $W$ be a word and $A$ a
finite set of additional letters such that $\D_A(W)=q$ and every letter of
$A$ is clumped in $W$. Then
\begin{equation}\label{eq:main}
 T(W)=r\quad\Longleftrightarrow\quad T(q)=r.
\end{equation}
There is no bound on $|A|$ or on the positive multiplicities of its letters.
The set $A$ may be empty. Adjacent equal letters in $q$, and insertions
between such letters, are allowed.
\end{theorem}

The assertion controls maximal time, not all sorting times.
The contribution is an insertion principle for words whose total alphabet
may be arbitrarily larger than that of $q$, together with the deletion
application in Section~\ref{sec:application}. The neighboring length counts
alone do not address this insertion. No equivalent statement was located
in the primary sources compared above in a search through 12 September
2026; this is a bounded literature comparison, not a claim of priority.
The original target motivating the application is to classify and count,
for every $k\geq3$, the length-$2k+2$ words on $k$ letters with $T=k$.
That full target remains unresolved in this note.

\section{Clumping and deletion}

We first recall the elementary structure behind the sorting bound. We
include the arguments to specify precisely the facts used in the insertion
proof.

\begin{lemma}[First-letter decomposition and clump growth]
\label{lem:clump}
If $w=a s_1 a s_2\cdots a s_m$, where $a$ occurs in none of the possibly
empty words $s_i$, then
\begin{equation}\label{eq:first}
 \phi(w)=\phi(s_1)\phi(s_2)\cdots\phi(s_m)a^m.
\end{equation}
Every clumped letter stays clumped. On each unsorted pass, the first
unclumped letter becomes clumped. Consequently a word with $u$ initially
unclumped letters has sorting time at most $u$.
\end{lemma}

\begin{proof}
The decomposition is Xia's first-letter lemma~\cite[Lemma 3.1]{Xia2024}.
The stack retains its bottom copies of $a$ while processing a segment $s_i$.
They cannot interact with the avoidance test inside that segment. Before
the next $a$ is pushed, every intervening non-$a$ entry must leave. Thus the
segments contribute their separate images in order; the copies of $a$
leave last. Final draining gives the same conclusion for $s_m$.

If a clumped block can begin to enter the stack, its remaining copies can
enter immediately. A later obstruction cannot be cleared by popping only
part of that block: any remaining copy gives the same obstruction. Hence
the block also leaves as a whole. Strip off leading clumped blocks of an
unsorted input and apply~\eqref{eq:first} to the remaining suffix. Its first
letter becomes clumped, and the stripped blocks move to the end without
splitting it. The number of unclumped letters therefore drops on every
unsorted pass. This is the clump-growth argument in
\cite[proof of Theorem 3.1]{Xia2024}, also stated in
\cite[Corollary 2.1]{ChoiGanLiZhu2025}.
\end{proof}

\begin{lemma}[Deletion of clumped letters]\label{lem:delete}
If $c$ is clumped in $w$, then
\begin{equation}\label{eq:delete}
 \D_c(\phi(w))=\phi(\D_c(w)).
\end{equation}
If all letters of $A$ are clumped in $w$, this identity extends to
$\D_A\phi^t(w)=\phi^t\D_A(w)$ for every $t\geq0$. In particular,
$T(\D_A(w))\leq T(w)$.
\end{lemma}

\begin{proof}
Induct on $|w|$. If $c$ is the first letter, write $w=c^m u$; both sides
of~\eqref{eq:delete} equal $\phi(u)$ by~\eqref{eq:first}. Otherwise, use
the first letter $a\ne c$ in~\eqref{eq:first}. The block of $c$ belongs to
exactly one segment $s_i$. Apply the induction hypothesis there and leave
the other segment images unchanged. This proves~\eqref{eq:delete}.
Clump persistence permits iteration. Deleting one clumped letter leaves
the other clumped letters clumped, so successive deletions give the assertion
for $A$. Finally, deleting letters from a sorted word preserves sortedness.
\end{proof}

For the pair $(q,W)$ in Theorem~\ref{thm:main}, these lemmas give
\begin{equation}\label{eq:bounds}
 T(q)\leq T(W)\leq r.
\end{equation}
The last bound uses the number of possibly unclumped old letters, rather
than the larger alphabet size of $W$.

\section{The finite base and induction}

Call the letters of $A$ \emph{markers}. Their allowed multiplicities and
number are unbounded, so the finite base requires the following reduction.

\begin{lemma}[Reduction to occupied gaps]\label{lem:gaps}
Fix an old word $q$ of length $n$. To test the sorting time of arbitrary
clumped-marker insertions into $q$, it suffices to insert one distinct
singleton marker in each gap belonging to a subset of its $n+1$ gaps.
More precisely, replacing all marker blocks in each occupied gap by one
fresh singleton preserves sortedness at every iterate.
\end{lemma}

\begin{proof}
First compress each marker block to one copy. Such a block is pushed and
popped as a unit, as in Lemma~\ref{lem:clump}. At the ends of its processing,
the stack transitions on old letters agree before and after compression.
Thus the old-letter output and the separations of old letters by markers
agree at every pass.

Within one original gap, the compressed markers form a string of distinct
singletons, none of whose letters occurs elsewhere. This string is pushed
without interruption. A later pop episode cannot stop strictly inside it:
the incoming letter differs from every member of the string, so an
obstruction requiring one of them to be popped also requires all remaining
members of that string above the obstruction to be popped. The whole string
is output consecutively in reverse order. The same property holds at the
next pass, since each member is still a singleton.

Replacing the string by a fresh singleton therefore preserves old-letter
transitions and whether markers separate two old runs. If marker strings
from different original gaps subsequently meet, neither can split. All
marker letters stay clumped in both systems, so agreement of the old-letter
order and separations implies agreement of full sortedness at every time.
Apply this replacement in every occupied gap. Marker names are immaterial
because relabelling commutes with $\phi$.
\end{proof}

\subsection*{Complete three-letter base}

Represent each word by renaming letters $0,1,2,\ldots$ in order of first
appearance. There are exactly
\[
 \frac{6!}{(2!)^3 3!}=15,
 \qquad
 \frac{7!}{3!(2!)^2 2!}=105
\]
representatives of profiles $(2,2,2)$ and $(3,2,2)$, respectively.
Every representative and every subset of its gaps were checked. In all
cases $\phi^2(W)$ is sorted if and only if $\phi^2(q)$ is sorted. The complete
counts are in Table~\ref{tab:base}; the empty gap subset is included.

\begin{table}[ht]
\centering
\small
\begin{tabular}{lrrrrr}
\toprule
Old profile & Old words & Masks/word & Marked inputs & $T(q)=3$ & $T(W)=3$\\
\midrule
$(2,2,2)$ & 15 & 128 & 1,920 & 1 & 128\\
$(3,2,2)$ & 105 & 256 & 26,880 & 12 & 3,072\\
\bottomrule
\end{tabular}
\caption{Exhaustive finite base. No sorting-equivalence failures occurred.}
\label{tab:base}
\end{table}

Appendix~\ref{app:code} contains a complete executable verifier using only
the Python standard library. It enumerates all labelled ternary words,
filters the multiplicities, and forms the canonical quotient. For each
proposed push, its stack routine checks every triple of positions for the
literal forbidden subsequence. Sortedness is checked from the intervals of
positions occupied by each letter. Both passes are recomputed on every
marked input; the code also checks deletion commutation in both passes.
Thus no sampling, numerical tolerance, or inferred recurrence is involved.

The checks also compare distinct exact formulations. In a related
combinatorial computation, Zeng's tree-reconstruction study computes
cut-partition data by integer dynamic programming and independently checks
small orders by direct edge-subset enumeration~\cite{Zeng2026MultiplicityArxiv}.
That study's Section~4 describes the checks as diagnostic, without using
them in the proof. Here the complete three-letter computation supplies
the induction base.

The original base calculation additionally compared a literal stack, a
run-stack implementation, and the recursion~\eqref{eq:first}. A separately
written verifier, using labelled words rather than that calculation's
multiplicity-pruned generator, reproduced the complete base; the appendix
is a compact presentation of this second computation. Since both sorting
times are at most three by~\eqref{eq:bounds}, testing the second iterate
proves~\eqref{eq:main} for the reduced three-letter inputs.
Lemma~\ref{lem:gaps} extends the base to arbitrary clumped markers.

\subsection*{Induction on the old alphabet}

\begin{proof}[Proof of Theorem~\ref{thm:main}]
The case $r=3$ was just established. Suppose $r\geq4$ and the theorem is
valid for $r-1$. Assume $T(W)=r$. By Lemma~\ref{lem:clump}, every old letter
is initially unclumped in $W$. Let $a$ be the first letter of $q$. Only
clumped markers can precede its first occurrence in $W$, so $a$ is the
first unclumped letter of $W$. It becomes clumped in $U=\phi(W)$; it is
also clumped in $v=\phi(q)$ by~\eqref{eq:first}.

Set $q_1=\D_a(v)$. Deletion commutation gives
\[
 \D_{A\cup\{a\}}(U)=q_1.
\]
All deleted letters are clumped in $U$. Moreover, $q_1$ has $r-1$ letters
with exactly their old multiplicities, so each occurs twice, except for at
most one third occurrence. It need not be run-free, which is why that
restriction was not imposed in the theorem. The induction hypothesis
therefore applies to $(q_1,U)$. Since $T(U)=r-1$, it yields $T(q_1)=r-1$.

By Lemma~\ref{lem:delete}, $T(v)\geq T(q_1)=r-1$. In particular $q$ cannot
be sorted, because its image and every deletion of that image would then
be sorted. Hence $T(q)=1+T(v)\geq r$. Together with~\eqref{eq:bounds} this
gives $T(q)=r$. Conversely, $T(q)=r$ and~\eqref{eq:bounds} force $T(W)=r$.
This completes the induction without invoking a two-letter case.
\end{proof}

\section{First-letter deletion at the next length}\label{sec:application}

A word is \emph{run-free} if adjacent letters are unequal. Its truncation
$\trunc(w)$ replaces each maximal constant run by one copy of its letter.

\begin{corollary}[Terminal insertion]\label{cor:terminal}
Let $q$ be run-free on $r\geq3$ letters, each occurring at least twice, and
let $2r\leq|q|\leq2r+1$. For a new letter $x$,
\[
 T(qx)=r\quad\Longleftrightarrow\quad T(q)=r.
\]
\end{corollary}

\begin{proof}
The length bounds give precisely the two profiles in
Theorem~\ref{thm:main}. Apply it with the single marker $x$.
\end{proof}

\begin{proposition}[First-letter deletion]\label{prop:application}
Let $k\geq4$, and let $w$ use exactly $k$ letters, with $|w|=2k+2$ and
$T(w)=k$. If its first letter $x$ occurs $m$ times, write
$\phi(w)=v x^m$ and put $r=k-1$ and $q=\trunc(v)$. Then $q$ uses exactly
$r$ letters, each at least twice, and
\[
 2r\leq |q|\leq2r+1,\qquad T(q)=r.
\]
\end{proposition}

\begin{proof}
Maximal time and Lemma~\ref{lem:clump} force zero clumped letters initially
and exactly one after the first pass, as in
\cite[Corollary 2.2]{ChoiGanLiZhu2025}. Thus every letter initially occurs
at least twice, $m\in\{2,3,4\}$, and the other $r$ letters are all
unclumped in $v$. Each therefore occupies at least two runs, giving the
lower bound $|q|\geq2r$.

If $m\geq3$, then $|v|=2k+2-m\leq2r+1$. If $m=2$, some other letter
occurs at least three times. Write $w=x s_1 x s_2$. That letter repeats
inside at least one $s_i$. The image of any word containing a repetition
has an adjacent equal pair: if its first letter repeats,
\eqref{eq:first} produces a terminal repeated block; otherwise remove its
unique first letter and apply the same argument to the shorter suffix.
Consequently $v=\phi(s_1)\phi(s_2)$ contains an adjacent equal pair. Its
truncation loses at least one position, giving
$|q|\leq |v|-1=2r+1$ also in this case.

Truncation compatibility is
\begin{equation}\label{eq:trunc}
 \trunc(\phi(z))=\trunc(\phi(\trunc(z))),
\end{equation}
as proved in~\cite[Proposition 3.1]{ChoiGanLiZhu2025}. To see the stack
mechanism, copies in an input run enter consecutively and form one stack
block; when that block must leave, retaining any positive number of its
copies leaves the same obstruction. Compressing runs therefore preserves
the sequence of output runs. Iterating~\eqref{eq:trunc} shows that $z$ and
$\trunc(z)$ have equal truncated iterates. Sortedness is equivalent to
having no repeated letter in the truncation, so
$T(z)=T(\trunc(z))$.

Because $x$ does not occur in $v$, we have $\trunc(vx^m)=qx$. It follows
that
\[
 T(qx)=T(vx^m)=k-1=r.
\]
Corollary~\ref{cor:terminal} now gives $T(q)=r$. The assumption $k\geq4$
ensures $r\geq3$ throughout.
\end{proof}

This reduction permits the first terminal clumped letter to be deleted
while retaining maximal time in one of the two preceding length layers.
It does not classify all inverse images of the resulting words or give
the requested count $M_k$ of length-$2k+2$ classes.

\section{Boundaries and verification scope}

Three exact examples distinguish the hypotheses and conclusion:
\[
\begin{array}{c|c|c|c}
 q & \text{terminal marker }x & T(q) & T(qx)\\ \hline
 0101 & 2 & 1 & 2\\
 0102102 & 3 & 1 & 2\\
 01201202 & 3 & 2 & 3
\end{array}
\]
The first excludes $r=2$. The second has the allowed profile $(3,2,2)$,
but shows that submaximal times need not be equal. The third has two
tripled letters and shows that arbitrary repeated-letter profiles do not
have the asserted maximal-time invariance. These values follow by direct
iteration of the displayed stack algorithm.

The proof uses only the complete three-letter base and the unbounded
reduction and induction above. Additional computations on $58{,}905$ marked
inputs with four or five old letters, and a separate $480$-input check of
expanded marker blocks in the three-letter base, gave no failures. They
are corroboration, not substitutes for the infinite argument. No all-$k$
enumeration is inferred from finite counts.

\paragraph{Computational and AI assistance.}
AI assistance was used in developing the argument, generating and checking
programs, comparing literature, and preparing this manuscript. A separate
AI-assisted review reconstructed the proof and reran the finite base with
a separately written standard-library verifier. Separate model contexts
are not external peer review. The initial literal simulation used Numba
to compile integer loops and was compared with two ordinary Python
implementations. The verifier printed below needs neither Numba nor any
external package. The mathematical decisions use equality tests and exact
integer arithmetic; the general theorem also depends on the written
compression and induction proofs. No formal proof-assistant certification
is claimed.

\appendix
\section{Executable finite-base verification}\label{app:code}

The following complete Python 3 program is supplied as
\texttt{verify\_base.py}. Run \texttt{python3 verify\_base.py}. It prints
the table rows and all $120$ old representatives with their second-iterate
sortedness and number of masks. In each row, the entries are length,
number of old representatives, number of marked inputs, number of old
words with time three, and number of marked inputs with time three.
Failure raises an exception; the checks are not disabled by Python's
optimization option. The generation visits $3^6+3^7$ labelled old words
before taking the quotient, and the resulting marker domain is exactly
the one in Lemma~\ref{lem:gaps}.

{\footnotesize
\begin{verbatim}
"""Exact finite base; run with Python 3. No external packages needed."""
from collections import Counter
from itertools import combinations, product
import json

def require(ok):
    if not ok:
        raise RuntimeError("finite-base verification failed")

def canonical(w):
    names = {}
    return tuple(names.setdefault(x, len(names)) for x in w)

def forbidden(s):
    return any(s[i] == s[k] != s[j]
               for i, j, k in combinations(range(len(s)), 3))

def phi(w):
    stack, out = [], []  # The top of the stack has index zero.
    for x in w:
        while forbidden([x] + stack):
            out.append(stack.pop(0))
        stack.insert(0, x)
    return tuple(out + stack)
\end{verbatim}
}
\newpage
{\footnotesize
\begin{verbatim}
def sorted_word(w):
    for x in set(w):
        p = [i for i, y in enumerate(w) if y == x]
        if p[-1] - p[0] + 1 != len(p):
            return False
    return True

def marked(q, mask):
    w = []
    for gap in range(len(q) + 1):
        if mask & (1 << gap):
            w.append(3 + gap)
        if gap < len(q):
            w.append(q[gap])
    return tuple(w)

rows, cases = [], []
for n, size in ((6, 15), (7, 105)):
    words = sorted({canonical(w) for w in product(range(3), repeat=n)
                    if sorted(Counter(w).values()) == [2, 2, n-4]})
    require(len(words) == size)
    hard_old = hard_marked = total = 0
    for q in words:
        old_second = phi(phi(q))
        old_hard = not sorted_word(old_second)
        hard_old += old_hard
        for mask in range(1 << (n+1)):
            u, v = marked(q, mask), q
            for step in range(2):
                u, v = phi(u), phi(v)
                require(tuple(x for x in u if x < 3) == v)
            require(sorted_word(u) == sorted_word(old_second))
            hard_marked += not sorted_word(u)
            total += 1
        cases.append({"q": q, "sorted_after_two": not old_hard,
                      "masks": 1 << (n+1)})
    rows.append([n, size, total, hard_old, hard_marked])
require(rows == [[6, 15, 1920, 1, 128], [7, 105, 26880, 12, 3072]])
print(json.dumps({"rows": rows, "cases": cases}, indent=2))
\end{verbatim}
}

\bibliographystyle{plainurl}
\bibliography{references}
\end{document}
